\documentclass[tikz,border=6pt]{standalone}
\usepackage[no-math]{fontspec}
\usepackage[brazil]{babel}
\usepackage{amsmath,amssymb}
\IfFileExists{C:/Windows/Fonts/Lato_Regular.ttf}{%
  \setmainfont{Lato}[
    Path = C:/Windows/Fonts/,
    UprightFont = Lato_Regular.ttf,
    BoldFont = Lato_Bold.ttf,
    ItalicFont = Lato_Italic.ttf,
    BoldItalicFont = Lato_BoldItalic.ttf
  ]
  \setsansfont{Lato}[
    Path = C:/Windows/Fonts/,
    UprightFont = Lato_Regular.ttf,
    BoldFont = Lato_Bold.ttf,
    ItalicFont = Lato_Italic.ttf,
    BoldItalicFont = Lato_BoldItalic.ttf
  ]
}{%
  \RequirePackage[default]{lato}
}
\renewcommand{\familydefault}{\sfdefault}

\usepackage{xcolor}
\usepackage{tikz}
\usepackage{pgfplots}
\pgfplotsset{compat=1.18}
\usetikzlibrary{arrows.meta,positioning,shapes.geometric,fit,backgrounds,calc,decorations.pathreplacing,patterns}

% Paleta institucional UEFS / CHF662
\definecolor{uefsNavy}{HTML}{1B2A4A}
\definecolor{uefsTeal}{HTML}{0E7C7B}
\definecolor{uefsForest}{HTML}{2D6A4F}
\definecolor{uefsGold}{HTML}{C4983B}
\definecolor{uefsGoldDark}{HTML}{8F6B1E}
\definecolor{uefsAlert}{HTML}{B02A37}
\definecolor{uefsOrange}{HTML}{D95F02}
\definecolor{uefsInk}{HTML}{22252A}
\definecolor{uefsMuted}{HTML}{5C6470}
\definecolor{uefsWater}{HTML}{2E86AB}
\definecolor{uefsSand}{HTML}{D9B382}
\definecolor{bgNavy}{HTML}{F2F5F9}
\definecolor{bgTeal}{HTML}{E9F5F4}
\definecolor{bgForest}{HTML}{EEF5F0}
\definecolor{bgGold}{HTML}{FCF7EB}
\definecolor{bgAlert}{HTML}{FDF2F2}
\definecolor{bgGray}{HTML}{F8FAFC}
\definecolor{lineGray}{HTML}{CFD8DC}

\tikzset{
  >=Stealth,
  aaCard/.style={draw=#1, fill=white, rounded corners=4pt, line width=1pt, inner sep=6pt, align=center, font=\small},
  aaBox/.style={draw=#1, rounded corners=4pt, line width=1.1pt, inner sep=6pt, align=center, font=\footnotesize\sffamily, text=uefsInk},
  aaTitle/.style={font=\normalsize\bfseries\sffamily, text=uefsNavy, align=center},
  aaSub/.style={font=\footnotesize\sffamily, text=uefsMuted, align=center},
  aaArrow/.style={->, >=Stealth, line width=1.4pt},
  aaDArrow/.style={<->, >=Stealth, line width=1.4pt}
}


\begin{document}
\begin{tikzpicture}

% =========================================================================
% Painel (a): Sistema Ambiental Aberto e Retroalimentação
% =========================================================================
\begin{scope}
  \node[aaTitle] at (5.1, 4.95) {(a) Sistema Ambiental Aberto e Retroalimentação};

  % Entradas
  \node[aaBox=uefsTeal, fill=bgTeal, text width=2.25cm, minimum height=2.05cm] (ent) at (1.05, 2.95)
    {\textbf{\color{uefsTeal}Entradas}\\[2pt]
     {\scriptsize\textbf{(Inputs)}}\\[4pt]
     {\scriptsize Chuva, radiação\\solar, insumos}};

  % Sistema Geoambiental (caixa externa)
  \node[draw=uefsNavy, fill=bgNavy, rounded corners=5pt, line width=1.2pt,
        minimum width=4.8cm, minimum height=2.55cm] (sis) at (5.1, 2.95) {};
  \node[font=\footnotesize\bfseries\sffamily, text=uefsNavy] at (5.1, 3.92) {Sistema Geoambiental};

  % Subcaixas Estoques e Processos
  \node[aaBox=uefsForest, fill=white, text width=1.72cm, minimum height=1.45cm] (est) at (3.95, 2.68)
    {\textbf{\color{uefsForest}Estoques}\\[3pt]
     {\scriptsize solo, água,\\biomassa}};

  \node[aaBox=uefsGoldDark, fill=white, text width=1.72cm, minimum height=1.45cm] (proc) at (6.25, 2.68)
    {\textbf{\color{uefsGoldDark}Processos}\\[3pt]
     {\scriptsize infiltração,\\erosão, uso}};

  \draw[aaDArrow, color=uefsNavy, line width=1.1pt] (est.east) -- (proc.west);

  % Saídas
  \node[aaBox=uefsAlert, fill=bgAlert, text width=2.25cm, minimum height=2.05cm] (sai) at (9.15, 2.95)
    {\textbf{\color{uefsAlert}Saídas}\\[2pt]
     {\scriptsize\textbf{(Outputs)}}\\[4pt]
     {\scriptsize Vazão, sedimento,\\calor, produção}};

  % Fluxo principal
  \draw[aaArrow, color=uefsTeal] (ent.east) -- (sis.west);
  \draw[aaArrow, color=uefsAlert] (sis.east) -- (sai.west);

  % Seta curva inferior de Retroalimentação passando POR BAIXO da caixa do Sistema Geoambiental
  \draw[aaArrow, dashed, color=uefsGoldDark, line width=1.35pt]
    (sai.south) .. controls (8.7, 0.72) and (1.5, 0.72) .. (ent.south);

  % Legenda da retroalimentação abaixo da curva
  \node[aaBox=uefsGold, fill=bgGold, inner sep=4.5pt, font=\scriptsize\sffamily, text=uefsGoldDark]
    at (5.1, 0.42) {\textbf{Retroalimentação (Feedback):} negativa = estabiliza; positiva = amplifica degradação};
\end{scope}

% =========================================================================
% Painel (b): Acoplamento Escala Espacial x Temporal
% =========================================================================
\begin{scope}[xshift=12.4cm, yshift=-0.05cm]
  \begin{axis}[
    width=7.5cm,
    height=6.2cm,
    title={(b) Acoplamento Escala Espacial $\times$ Temporal},
    title style={font=\normalsize\bfseries\sffamily, text=uefsNavy, yshift=2pt},
    axis lines=left,
    axis line style={lineGray, line width=1pt},
    tick style={uefsMuted},
    xmin=-0.35, xmax=5.65,
    ymin=-0.15, ymax=5.05,
    xtick={0.5, 2.7, 4.7},
    xticklabels={{Local\\(parcela)}, {Meso\\(bacia)}, {Macro\\(regional)}},
    xticklabel style={align=center, font=\scriptsize\sffamily, text=uefsInk},
    ytick={0.5, 2.4, 4.2},
    yticklabels={Evento, Anos, Séculos},
    yticklabel style={font=\scriptsize\sffamily, text=uefsInk},
    xlabel={\textbf{Escala Espacial}},
    ylabel={\textbf{Escala Temporal}},
    label style={font=\footnotesize\sffamily, text=uefsInk},
    clip=true
  ]
    % Linha de tendência
    \addplot[dashed, color=lineGray, line width=1.1pt] coordinates {(0.5,0.5) (4.7,4.2)};

    % Pontos individuais
    \addplot[only marks, mark=*, mark size=4.2pt, color=uefsAlert] coordinates {(0.5,0.5)};
    \addplot[only marks, mark=*, mark size=4.2pt, color=uefsGoldDark] coordinates {(1.6,1.5)};
    \addplot[only marks, mark=*, mark size=4.2pt, color=uefsTeal] coordinates {(2.7,2.4)};
    \addplot[only marks, mark=*, mark size=4.2pt, color=uefsNavy] coordinates {(3.8,3.3)};
    \addplot[only marks, mark=*, mark size=4.2pt, color=uefsForest] coordinates {(4.7,4.2)};

    % Rótulos deslocados sem sobrepor pontos nem bordas
    \node[anchor=west, align=left, font=\scriptsize\bfseries\sffamily, text=uefsAlert]
      at (axis cs:0.72, 0.42) {Splash / Selamento\\[-1pt]{\tiny (cm\textsuperscript{2} | minutos)}};

    \node[anchor=north west, align=left, font=\scriptsize\bfseries\sffamily, text=uefsGoldDark]
      at (axis cs:1.78, 1.38) {Sulcos e Ravinas\\[-1pt]{\tiny (vertente | chuvas)}};

    \node[anchor=north west, align=left, font=\scriptsize\bfseries\sffamily, text=uefsTeal]
      at (axis cs:2.85, 2.28) {Microbacia / Voçoroca\\[-1pt]{\tiny (km\textsuperscript{2} | anos)}};

    \node[anchor=south east, align=right, font=\scriptsize\bfseries\sffamily, text=uefsNavy]
      at (axis cs:3.65, 3.45) {Bacia Regional\\[-1pt]{\tiny (10\textsuperscript{3} km\textsuperscript{2} | décadas)}};

    \node[anchor=south east, align=right, font=\scriptsize\bfseries\sffamily, text=uefsForest]
      at (axis cs:4.55, 4.30) {Bioma / Clima\\[-1pt]{\tiny (10\textsuperscript{6} km\textsuperscript{2} | séculos)}};
  \end{axis}
\end{scope}

\end{tikzpicture}
\end{document}
